\section*{ОБОЗНАЧЕНИЯ И СОКРАЩЕНИЯ}

КА -- конечный автомат.

ДКА -- детерминированный конечный автомат.

НКА -- недетерминированный конечный автомат.

АСД -- абстрактное синтаксическое дерево.

BNF (Backus Naur Form) -- формальная грамматика в форме Бэкуса-Наура.

API (Application Programming Interface) -- набор правил, типов, имен функций и контрактов, через которые одна программная сущность взаимодействует с другой на уровне исходного кода.

ABI (Application Binary Interface) -- соглашение о взаимодействии на уровне двоичного кода (машинный код, объектные файлы, библиотеки) между компилированными компонентами.

UB (Undefined Behavior) -- неопределенное поведение; ситуация, когда в определенных граничных случаях поведение программного продукта может меняться неконтролируемым образом и приводить к некорректным результатам.

ПО -- программное обеспечение.

РП -- рабочий проект.

ТЗ -- техническое задание.

ТП -- технический проект.

UML (Unified Modelling Language) -- язык графического описания для объектного моделирования в области разработки программного обеспечения.
